\section{Time frequency transform with frequency focus}\label{sec_frequency_focus}
\label{sec:freqfocus}

\subsection{Definition of the transform}

We now introduce the frequency focus effect, generated by the associated frequency focus function $\sigma^\nu$. The role of the focus function is to modify the shape of the analysis waveforms, in a way that depends on some local behaviour of the analysed signal $f$. Starting from the construction described in Section~\ref{sec_CSQT}, the time-frequency atoms $h_{t,\u}$ are locally rescaled using a so called focus function $\sigma^\nu_f$, which depends on $f$. This yields new atoms that will be denoted by $h_{t,\u,\sigma^\nu_f}$, $t$ denoting time and $u$ being the frequency scaling .

\paragraph{Assumptions} As in the previous section, $\gamma$ is assumed to be a positive, increasing $C^1$ diffeomorphism that maps $\R$ onto $\R_+$. The focus functions $\sigma^\nu_f$ of $f\in L^2(\R)$ will be assumed to be at least piecewise continuous, larger than 1 and such that $\sigma_f^\nu -1$ is compactly supported. To sum up
\begin{equation}
	\forall f\in L^2(\R), \ \sigma^\nu_f \in C_{p,c}^+(\R) + 1\ .
\end{equation}



\paragraph{Atoms and transform} Given these notations, the Viennese constant-Q transform with frequency focus effect (or frequency-Marseillan transform) $M^\nu$ can be defined for $f\in C_c$ by

\begin{equation}
	\forall t,\u\in\R\ ,\quad \ M^\nu f(t,\u) := \langle f , h_{t,\u,\sigma^\nu_f} \rangle_{L^2_x}\ ,
\end{equation}
the time-frequency atoms being defined by
\begin{equation}
	\forall x,t,\u\in\R\, , \quad h_{t,\u,\sigma_f^\nu}(x) := \sigma_f^\nu(\u)\sqrt{\gamma(\u)}e^{2i\pi\gamma(\u)x}h((\gamma(\u)x-t)\sigma_f^\nu(\u))\ .
\end{equation}

\begin{remark}
In order to keep constant $L^2$ norm for the atoms, we should normalize them by $\sqrt{\sigma_f}$ instead of $\sigma_f$, however with such normalisation $d\mu$ will not be a measure anymore which goes beyond the borders of our study.
\end{remark}

The main purpose of the section is to establish that $M^\nu$ can be extended to a well defined map from $L^2(\R,dx)\rightarrow L^2(\R^2,d\mu)$ satisfying a norm control similar to a frame bound control.  The pointwise definition of $M^\nu f(t,\u)$ on $L^2$ is still guaranteed by the fact that $h\in L^2$. Regarding the definition from $L^2$ into $L^2(\R^2,d\mu)$ we can raise that $M^\nu$ is well defined from $C_c$ into $L^2$ and then extend the control by the use of the Fatou's lemma.

By computations similar as the previous section, we obtain the following Fourier transform of the atoms,
\begin{equation}
	\F(h_{t,\u,\sigma^\nu(\u)})(\xi) = \frac{1}{\sqrt{\gamma(\u)}}\hat{h}\left( \beta_\xi(\u)/\sigma^\nu(\u) \right)e^{-2i\pi\beta_\xi(\u)t}\ .
\end{equation}
Sometimes, if there are no ambiguities we will write $h_{t,\u,f}$ instead of $h_{t,\u,\sigma_f}$ and $h_{t,\u,n}$ instead of $h_{t,\u,\sigma_{f_n}}$ for a certain sequence $(f_n)_n$.

\subsection{Example and ideas of the frequency focus function}\label{ssec_frequency_focus_function}
Defining and studying specific frequency focus functions is not a main goal of this paper. We can however provide a toy example that could be of interest in some specific situations.
\begin{example}[Resolving close frequencies]
\label{ex:frequency_focus}
	Imagine a function $f$ of interest contains several harmonic components (i.e. sine waves) with more or less close frequencies. When consecutive frequencies are too close, a wideband window will not be able to separate them, which would suggest to increase the frequency precision in the neighborhood of these frequencies. A prototypical choice for frequency focus function in such a context could be
	\begin{equation}
 \label{fo:ex.freq.foc}
		\sigma_f^\nu(\u) = 1 + \frac{\|V_gf(\cdot,\u)\|_{L^1(dt)}}{\|V_gf(\cdot,\u)\|_{L^2(dt)}}\,\varphi(\u)\ ,
	\end{equation}
	where $\varphi$ is compactly supported in some interval $[-\u_0,\u_0]$ for some cutoff frequency $\u_0$ ($\varphi$ can be a multiple of the indicator function of that interval, or the hearing threshold in psycho-physical applications), and $V_g$ is the ordinary STFT, with gaussian window~\cite{grochenig2013foundations}.\\
	%
	The ratio of $L^1$ to $L^2$ norms is known as a a measure of spread, or diversity (it takes small values when the corresponding function is more concentrated, and is closely related to a specific Renyi entropy, see~\cite{Ricaud2013refined}). When $f$ contains several sine waves in the neighbourhood of $\u$, which are not separated by the gaussian window, some beats appear in the function $t\mapsto V_gf(t,\u)$, which yields smaller $L^1/L^2$ norm values, and therefore small values for $\sigma_f^\nu(\u)$.

\smallskip
Illustrations with a suitably post-processed version of this choice, on a toy signal example, is provided in Appendix~\ref{sec:app.freq.foc.example}.
%

\end{example}
\subsection{Kernel and norm relationship}

To prove the main result we first derive a norm relationship involving a certain non-negative valued kernel $K_\sigma$, so that the study of the norm of $M^\nu f$ will be determined by the norm $\|K_\sigma\|_\infty$.
\begin{prop}
	Let us assume that $h$ satisfies the same admissibility condition as in  Proposition~\ref{prop_norm_relation_operator_Q_transform}. Then we have
	\begin{equation}
		\|M^\nu f\|_{L^2(d\mu)}^2 = \int_\R |\hat{f}(\xi)|^2 K_{\sigma_f^\nu}(\xi)d\xi\ ,
	\end{equation}
	where the kernel $K_{\sigma_f^\nu}(\xi)$ is defined by 
	\begin{equation}\label{kernel_frequency_case}
		K_{\sigma_f^\nu}(\xi) := \int_{-1}^{+\infty} \left|\hat{h}\left( \frac{y}{\sigma_f^\nu\circ\gamma^{-1}(|\frac{\xi}{y+1}|)} \right)\right|^2\frac{dy}{|y+1|}\ .
	\end{equation}
\end{prop}
\begin{remark}
Since $\sigma\circ\gamma^{-1}(|\frac{\xi}{y+1}|)=1$ for $y$ big enough or close enough to $-1$, thanks to the admissiblity condition on $h$  (see equation \ref{CQT.admissibility}) the convergence of $K_\sigma$ is guaranteed.
\end{remark}
\begin{proof}
	Introducing $\sigma(\u)$ in the proof of Proposition~\ref{prop_norm_relation_operator_Q_transform} gives us 
	\begin{align*}
		\|M^\nu f\|^2_{L^2} & = \int \bar{\hat{f}}(\xi)\hat{f}(\xi')\frac{1}{\gamma(\u)}\hat{h}(\beta_\xi(\u)/\sigma_f^\nu(\u))\bar{\hat{h}}(\beta_{\xi'}(\u)/\sigma_f^\nu(\u))e^{-2i\pi\gamma(\u)^{-1}(\xi-\xi')t}d\xi d\xi'd\mu(\u,t) \\
		& = \int |\hat{f}(\xi)|^2 \frac{\gamma'(\u)}{\gamma(\u)^2} |\hat{h}(\beta_\xi(\u)/\sigma_f^\nu(\u)|^2 d\xi d\u \\
		& = (1) + (2)\ ,
	\end{align*}
	where
	\begin{align*}
		(1) &:= \int_0^{+\infty}|\hat{f}(\xi)|^2\int_{-\infty}^{+\infty}\frac{\gamma'(\u)}{\gamma(\u)^2} |\hat{h}(\beta_\xi(\u)/\sigma_f^\nu(\u)|^2 d\u d\xi, \\
		(2) & := \int_{-\infty}^0|\hat{f}(\xi)|^2\int_{-\infty}^{+\infty}\frac{\gamma'(\u)}{\gamma(\u)^2} |\hat{h}(\beta_\xi(\u)/\sigma_f^\nu(\u)|^2 d\u d\xi\ .
	\end{align*}
	Using the same change of variable $y= \beta_\xi(\u)$ in both cases and taking care of the signs we get (1) and in (2) the sign of $\xi/(y+1)$ is still non negative as $\xi\leqslant 0$ and $y\leqslant -1$ so we obtain
	\begin{align*}
		(1) &=  \int_0^{+\infty}|\hat{f}(\xi)|^2\int_{-1}^{+\infty} \left|\hat{h}\left( \frac{y}{\sigma_f^\nu\circ\gamma^{-1}(\frac{\xi}{y+1})} \right)\right|^2\frac{dy}{y+1} d\xi , \\
		(2) &= -\int_{-\infty}^0|\hat{f}(\xi)|^2 \int_{-\infty}^{-1} \left|\hat{h}\left( \frac{y}{\sigma_f^\nu\circ\gamma^{-1}(\frac{\xi}{y+1})} \right)\right|^2\frac{dy}{y+1} d\xi \\
		&= \int_{-\infty}^0|\hat{f}(\xi)|^2 \int_{-\infty}^{-1} \left|\hat{h}\left( \frac{y}{\sigma_f^\nu\circ\gamma^{-1}(|\frac{\xi}{y+1}|)} \right)\right|^2\frac{dy}{|y+1|} d\xi.
	\end{align*}
	Using a symmetry argument similar to the proof of proposition~\eqref{prop_norm_relation_operator_Q_transform}, and summing both terms we obtain the expected formula.
\end{proof}

	\begin{remark}
		If $\sigma$ is fixed and independent of $f$, the time-frequency atoms $h_{t,\omega,\sigma(\omega)}$ form a \textit{continuous frame} of $L^2(\R)$ in the terminology of~\cite{Ali2000coherent,Dahlke2008generalized}. We will see that in the general case, the assumptions made on $\sigma$ allow one to stay in a tractable situation.
\end{remark} 
From now, the goal is to obtain upper and lower bounds for the kernel $K_\sigma$.

\subsection{Main result : the $L^2$ norm control}\label{ssection_norm_control_frequency_focus}

The main result of this section is the following theorem~\ref{theorem_norm_control_frequency_focus} which is a generalisation to non linear transform with adaptive window of the classical frame control.
\begin{theorem}\label{theorem_norm_control_frequency_focus}
	Let $f\in L^2(\R)$ and $\sigma^\nu_f\in C_c^+(\R)\cap C_p+1$ with support included in $[a,b]$, $a\leqslant b$. Then we have
	\begin{equation}
		c_{a,b}\|f\|_{L^2(dx)} \leqslant \|M^\nu f \|_{L^2(d\mu)} \leqslant C_{a,b}\|f\|_{L^2(dx)},
	\end{equation}
	where the constants $c_{a,b}$ and $C_{a,b}$ are given by
	\begin{equation}
		C_{a,b}^2 := \int_{-1}^{+\infty}|\hat{h}(y)|^2 \frac{dy}{y+1} + \|\hat{h}\|_{\infty}^2\ln\left(\frac{\gamma(b)}{\gamma(a)}\right) < + \infty,
	\end{equation}
	and
	\begin{equation}
		c_{a,b}^2 := \inf_{\xi_0\in\R}\{r^-(\xi_0) + r^+(\xi_0) \} >0,
	\end{equation}
	with
	\begin{equation}
		r^-(\xi):= \int_{-1}^{-1+|\xi|/\gamma(b)}\frac{|\hat{h}(y)|^2}{y+1}dy \  , \quad r^+(\xi):= \int_{-1+|\xi|/\gamma(a)}^{+\infty}\frac{|\hat{h}(y)|^2}{y+1}dy\ .
	\end{equation}
\end{theorem}

\begin{remark}
	We can raise that if $\sigma$ is fixed (\ie\, independent of $f$) then we obtain the classical bound control from linear time frequency operators since $c_{a,b}$ and $C_{a,b}$ only depend on $\sigma$ and not directly on $f$. More precisely they only depend on the support of $\sigma$.
\end{remark}

\begin{proof}
	Even though $\sigma_f$ is supposed to be piecewise continuous, the idea of the proof is to control the sup norm of the kernel $K_\kappa$ when $\kappa \in \mathcal{E}_c^+ +1$ with $\mathcal{E}$ the set of step functions and $\kappa$ close enough to $\sigma$. Everything will be defined correctly in the next section. The reason of such idea is because for step functions $\kappa$ the kernel $K_\kappa$ is easy to estimate. Thus by using the control for step functions and proving a continuity result of $K_\kappa$ for functions $\kappa\rightarrow \sigma$ we will be able to prove our control for continuous (pr piecewise continuous) focus functions $\sigma_f$.
\end{proof}

\subsection{Proof of the theorem \ref{theorem_norm_control_frequency_focus} : approach by step functions }\label{ssection_rectangular_function_frequency_focus}

\subsubsection{Decomposition lemma}

We introduce the decomposition lemma \ref{lemma_decomposition_kernel_square_function} that is a key part of the study of the kernel onto step functions. Indeed it turns out that when $\kappa\in\mathcal{E}^+([a,b])+1$, then $K_\kappa$ can be expressed in quite a simple way which yields 'easy' estimates depending only on the support $[a,b]$ and $h$.

\begin{lemma}[Decomposition of the Kernel]
	\label{lemma_decomposition_kernel_square_function}
	Let $\kappa(x):=1+\lambda\mathds{1}_{[a,b]}(x)$, $a<b$ and $\lambda\geqslant0$, thus
	\begin{equation}
		\forall \xi\in\R,\qquad K_\kappa(\xi) = \int_{[-1,+\infty)\backslash [\beta_{|\xi|}(b),\beta_{|\xi|}(a)]}\frac{|\hat{h}(y)|^2}{|y+1|}\,dy\ +\  \int_{\beta_{|\xi|}(b)}^{\beta_{|\xi|}(a)}\frac{|\hat{h}(\tfrac{y}{1+\lambda})|^2}{|y+1|}\,dy\ .
	\end{equation}
	
\end{lemma}
\begin{proof}
	If $\xi=0$ then the lemma is clearly true, so let us assume $\xi\not=0$. Now let us raise that $\sigma\circ\gamma^{-1}(\u)=1+\mathds{1}_{[\gamma(a),\gamma(b)]}(\u)$ thus $\sigma\circ\gamma^{-1}(|\tfrac{\xi}{y+1}|)=1+\mathds{1}_{[\gamma(a),\gamma(b)]}(|\tfrac{\xi}{y+1}|)$. In order to have $\sigma\circ\gamma^{-1}(|\tfrac{\xi}{y+1}|)=1+\lambda$ there are two possible situations~: either $-1+\frac{\xi}{\gamma(b)} \leqslant y \leqslant -1+\frac{\xi}{\gamma(a)}$ (case (1)) or $-1-\frac{\xi}{\gamma(a)} \leqslant y \leqslant -1-\frac{\xi}{\gamma(b)}$ (case (2)). Since $y\geqslant -1$ only one case is possible at the time. When $\xi>0$ case (2) is not possible but (1) is, so we have the desired decomposition
	\begin{equation}
		K_\kappa(\xi) = \int_{[-1,+\infty)\backslash [\beta_\xi(b),\beta_\xi(a)]}\frac{|\hat{h}(y)|^2}{|y+1|}\,dy + \int_{\beta_\xi(b)}^{\beta_\xi(a)}\frac{|\hat{h}(\tfrac{y}{1+\lambda})|^2}{|y+1|}\,dy\ .
	\end{equation}
	However if $\xi<0$ only case (2) is possible and thus the decomposition is
	\begin{equation}
		K_\kappa(\xi) = \int_{[-1,+\infty)\backslash [\beta_{-\xi}(b),\beta_{-\xi}(a)]}\frac{|\hat{h}(y)|^2}{|y+1|}\,dy + \int_{\beta_{-\xi}(b)}^{\beta_{-\xi}(a)}\frac{|\hat{h}(\tfrac{y}{1+\lambda})|^2}{|y+1|}\,dy\ .
	\end{equation}
	Combining both decompositions yields the result.
\end{proof}
Following the same ideas we can easily generalise the decomposition lemma to more complex step functions.

\begin{lemma}\label{lemma_decomposition_kernel_square_function_general}
	Lets $\kappa(x) := 1 + \sum_{i=0}^n \lambda_i \mathds{1}_{[a_i,b_i]}(x)$ with $\lambda_i\geqslant 0$ and $a_i\leqslant b_i \leqslant a_{i+1}$. Then we have
	\begin{equation}
		\forall \xi\in\R\, ,\quad \ K_\kappa(\xi) = \int_{[-1,+\infty)\backslash [\beta_{|\xi|}(b),\beta_{|\xi|}(a)]}\frac{|\hat{h}(y)|^2}{|y+1|}dy + \sum_{i=0}^n\int_{\beta_{|\xi|}(b_i)}^{\beta_{|\xi|}(a_i)}\frac{|\hat{h}(\tfrac{y}{1+\lambda_i})|^2}{|y+1|}dy\ .
	\end{equation}
\end{lemma}
Besides giving us an important decomposition of the integral it also gives us the symmetry properties of our object, and will then allow us to restrain to $\xi\geqslant0$ all computations which involve this decomposition (as long as they are symmetric on $\xi$ on their others parts).

\subsubsection{Upper bound control}

Before introducing the general upper bound result we first prove a weaker result that concerns function of the form : $ \kappa(\u) = 1 + \lambda\mathds{1}_{[a,b]}(\u)$.


\begin{prop}\label{prop_control_kernel_of_square_function}
	Let $a<b$ and $\kappa(\u)= 1 + \lambda \mathds{1}_{[a,b]}(\u)$, then we have
	\begin{equation}
		\forall \xi, \ K_\kappa (\xi) \leqslant \int_{-1}^{+\infty}\frac{|\hat{h}(y)|^2}{|y+1|}dy + \|\hat{h}\|_{L^\infty}^2\ln\left(\frac{\gamma(b)}{\gamma(a)}\right).
	\end{equation}
\end{prop}
\begin{proof}
	Using Lemma~\ref{lemma_decomposition_kernel_square_function} we obtain for $\xi\geqslant0$
	\begin{align}
		K_\kappa(\xi) &= \int_{[-1,+\infty)\backslash [-1+\frac{\xi}{\gamma(b)},-1+\frac{\xi}{\gamma(a)}]}\frac{|\hat{h}(y)|^2}{|y+1|}dy + \int_{-1+\frac{\xi}{\gamma(b)}}^{-1+\frac{\xi}{\gamma(a)}}\frac{|\hat{h}(\tfrac{y}{1+\lambda})|^2}{|y+1|}dy \\
		& \leqslant \int_{-1}^{+\infty}\frac{|\hat{h}(y)|^2}{|y+1|}dy + \|\hat{h}\|_{L^\infty}^2\int_{-1+\frac{\xi}{\gamma(b)}}^{-1+\frac{\xi}{\gamma(a)}}\frac{1}{|y+1|}dy \\
		&= \int_{-1}^{+\infty}\frac{|\hat{h}(y)|^2}{|y+1|}dy + \|\hat{h}\|_{L^\infty}^2\ln\left(\frac{\gamma(b)}{\gamma(a)}\right).
	\end{align}
	The same computation can be performed in the case $\xi<0$.
\end{proof}
The norm $\|\hat{h}\|_{L^\infty}$ is easy to control since it is below $\|h\|_{L^1}$. What is interesting with this first result is the fact that the control is independent of $\lambda$. The result is easy to generalise, by splitting the integral into as many terms as there are subdivisions. This gives Proposition~\ref{prop_control_norm_kernel_on_stair_function} below.

\begin{prop}\label{prop_control_norm_kernel_on_stair_function}
	Let $(a_i)_{0\leqslant i \leqslant n},(b_i)_{0\leqslant i \leqslant n}$ such that $a_{i+1}>b_i$, $(\lambda_i)_{0\leqslant i \leqslant n}\in (\R_+^*)^{n+1}$ and $\sigma(x)= 1 + \sum_{i=0}^n \lambda_i\mathds{1}_{[a_i,b_i]}(x)$ then
	\begin{equation}
		\forall \xi \in \R, \ K_{\sigma}(\xi) \leqslant \int_{-1}^{+\infty}\frac{|\hat{h}(y)|^2}{|y+1|}dy + \|\hat{h}\|_{L^\infty}^2\sum_{i=0}^n\ln\left(\frac{\gamma(b_i)}{\gamma(a_i)}\right)\ .
	\end{equation}
	A consequence of such control is that there exists a constant $C^+_{a,b}$ that depends on the window $h$ and on $a,b$ such that for any $\kappa\in \mathcal{E}^+([a,b])+1$ we have
	\begin{equation}
		\|K_\kappa\|_\infty \leqslant C_{a,b}^2 := \int_{-1}^{+\infty}\frac{|\hat{h}(y)|^2}{|y+1|}dy + \|\hat{h}\|_{L^\infty}^2\ln\left(\frac{\gamma(b)}{\gamma(a)}\right)\ .
	\end{equation}
\end{prop}
\begin{proof}
	The proof is the same as the previous one except that Lemma~\ref{lemma_decomposition_kernel_square_function_general} is used and the contribution of every sub-interval $[-1+\xi/\gamma(b_i),-1+\xi/\gamma(a_i)]$ has to be controlled.
\end{proof}

Even though it is not the best possible control, this bound is strong enough for our purpose in Theorem \ref{theorem_norm_control_frequency_focus}. A stronger control for a specific $\kappa$ could be the following one but since it depends too much on the values of $(a_i),(b_i),(\lambda_i)$ we will not use it later.


\begin{prop}\label{prop_norm_control_kernel_stair_function_slighter}
	Let $\kappa(x)=\sum_{i=0}^n \lambda_i \mathds{1}_{[a_i,b_i]}(x)$ with $\lambda_i\geqslant0$ and $a_{i+1}\geqslant b_i$, then
	\begin{equation}
		\forall \xi\in\R,\quad |K_\kappa (\xi) | \leqslant \int_{-1}^{+\infty} \frac{|\hat{h}(y)|^2}{y+1}dy + \max_{0 \leqslant i \leqslant n}\sup_{[B_i^\xi,A_i^\xi]}\{|\hat{h}(y)|^2 \} \sum_{j=0}^n \ln\left( \frac{\gamma(b_j)}{\gamma(a_j)} \right)\ ,
	\end{equation}
	where $A_i^\xi=\frac{-1+\xi / \gamma(a_i)}{1+\lambda_i}, B_i^\xi= \frac{-1+\xi / \gamma(b_i)}{1+\lambda_i}$.
\end{prop}
\begin{proof}
	We use the same decomposition as previously, thus letting $0\leqslant i \leqslant n$ we have
	\begin{align}
		\int_{-1+\xi / \gamma(b_i)}^{-1+\xi / \gamma(a_i)}\frac{|\hat{h}(\tfrac{y}{1+\lambda_i})|^2}{y+1}dy &\leqslant \sup_{[B_i^\xi,A_i^\xi]}\{|\hat{h}(y)|^2\}\ln\left( \frac{\gamma(b_i)}{\gamma(a_i)}  \right)
	\end{align}
	and taking the max over $0\leqslant i \leqslant n$, summing the log and keeping the rest $\int |\hat h(y)|^2/(y+1)dy$ we obtain the desired result.
\end{proof}

\subsubsection{Lower bound control}

The following result guarantees the existence of a positive lower bound that only depends on the support of the focus function.

\begin{prop}\label{prop_lower_bound_kernel_basic}
	Let $a<b\in\R$ and $\kappa\in \CalE([a,b],\R^+)+1$, for any $\xi\in\R$ we define 
	\begin{equation}
		r^-(\xi):= \int_{-1}^{-1+|\xi|/\gamma(b)}\frac{|\hat{h}(y)|^2}{y+1}dy
	\end{equation}
	and
	\begin{equation}
		r^+(\xi):= \int_{-1+|\xi|/\gamma(a)}^{+\infty}\frac{|\hat{h}(y)|^2}{y+1}dy\ .
	\end{equation}
	Then we have
	\begin{equation}
		\inf_\xi |K_\kappa(\xi)|\geqslant \inf_{\xi_0\in\R}\{r^-(\xi_0) + r^+(\xi_0) \}=:c_{a,b}^2\ .
	\end{equation}
	Furthermore, under our hypothesis $c_{a,b}>0$.
\end{prop}
\begin{proof}
	A separation trick similar to what was done in Lemma~\ref{lemma_decomposition_kernel_square_function} gives us, for any $\xi\in\R$,
	\begin{align*}
		K_\kappa(\xi) &=  \int_{-1}^{-1+|\xi|/\gamma(b)}\frac{|\hat{h}(y)|^2}{y+1}dy \\
		& + \int_{-1+|\xi|/\gamma(b)}^{-1+|\xi|/\gamma(a)}\left|\hat{h}\left( \frac{y}{\kappa\circ\gamma^{-1}(|\frac{\xi}{y+1}|)} \right)\right|^2\frac{dy}{|y+1|} \\
		& +\int_{-1+|\xi|/\gamma(a)}^{+\infty}\frac{|\hat{h}(y)|^2}{y+1}dy.
	\end{align*}
	Hence by taking the $\inf$ we obtain the expected control, furthermore by continuity of $r^{\pm}(\xi)$ and by positivity at every $\xi\in\R$ and at the limits at $\pm\infty$ of at least $r^+$ or $r^-$ we can conclude on the positivity of $c_{a,b}$.
\end{proof}


\subsubsection{Continuity and general norm control}\label{sssection_continuity_norm_control_frequency_focus}


We recall that a function $f:\R\rightarrow\C$ is uniformly continuous if
\begin{equation}\label{eq_uniform_continuity}
	\forall \eps>0,\, \exists \delta(\eps)>0,\, \forall x,y\in\R,\, |x-y|\leqslant \delta(\eps) \Rightarrow |f(x)-f(y)|\leqslant \eps\ .
\end{equation}
Of course there is no uniqueness of $\delta(\eps)$ but we can take anyone that satisfies equation~\eqref{eq_uniform_continuity}. We can now state the continuity result that is the key point of the whole section.

\begin{prop}[Continuity of the Kernel]\label{prop_continuity_kernel}
	Assume that $|\hat{h}|\in C_0(\R)$, then $|\hat{h}|^2$ is uniformly continuous with constant $\delta(\eps)>0$ for $\eps>0$. And let $\sigma\in C([a,b])^+\cap C_p +1$. Then there exists $\xi_0(\eps)>0$ such that for any $\kappa\in \mathcal{E}([a,b])+1$ satisfying $\kappa(x)\leqslant \sigma(x)$ and
	\begin{equation}
		\| \sigma - \kappa \|_{\infty} \leqslant \tilde{\delta}(\eps) := \frac{\delta(\eps)\gamma(a)}{\xi_0(\eps)}\ ,
	\end{equation}
	we have
	\begin{equation}
		\| K_\sigma - K_\kappa \|_{\infty} \leqslant \eps \ln\left(\frac{\gamma(b)}{\gamma(a)}\right)\ .
	\end{equation}
\end{prop}
\begin{proof}
	
	Let us first raise that for $\sigma\in C^+([a,b])+1$ and $\kappa\in\mathcal{E}^+([a,b])+1$ then we have for any $\xi\in\R$,
	\begin{equation*}
		|K_\sigma(\xi)-K_\kappa(\xi)| = \left|\int_{-1+\frac{|\xi|}{\gamma(b)}}^{-1+\frac{|\xi|}{\gamma(a)}} \left|\hat{h}\left(\frac{y}{\sigma\circ\gamma^{-1}(\frac{\xi}{y+1})}\right)\right|^2-\left|\hat{h}\left(\frac{y}{\kappa\circ\gamma^{-1}(\frac{\xi}{y+1})}\right)\right|^2 \frac{dy}{y+1} \right|\ .
	\end{equation*}
	Therefore, for a fixed $\xi$ we will be working on bounded intervals, which will be is important for the rest of the proof. The main idea is the following  : consider the quantity $|K_\kappa(\xi)-K_\sigma(\xi)|$ in two cases. The first one will be for $\xi\geqslant\xi_0$ large enough such that the integrated quantity will be small because $\hat{h}$ tends to infinity at $\pm\infty$. The second case will be for $\kappa$ close enough with respect to $\eps$ and $\xi_0$ so that the difference in the integral will be small enough.
	
	Since $|\hat{h}|$ tends to $0$ at $\pm\infty$ we have the existence of $y_0(\eps)$ such that for any $|y|\geqslant y_0$ the quantity $|\hat{h}(y)|^2 \leqslant \eps/2$. Hence there exist $\xi_0(\eps)$ such that $\xi_0/\gamma(b)\geqslant 1$ and $\frac{\xi_0/\gamma(b) -1}{\|\sigma\|_\infty} \geqslant y_0(\eps)$.
	Furthermore since for any $\xi\geqslant \xi_0$ and any $y\in[-1+\frac{\xi}{\gamma(b)},-1+\frac{\xi}{\gamma(a)}]$ we have
	\begin{equation}
		\frac{y}{\sigma\circ\gamma^{-1}(\frac{|\xi|}{y+1})} \geqslant \frac{\frac{|\xi|}{\gamma(b)}-1}{\|\sigma\|_\infty}\geqslant \frac{\frac{|\xi_0|}{\gamma(b)}-1}{\|\sigma\|_\infty} \geqslant y_0(\eps)\ ,
	\end{equation}
	then we have the following control
	\begin{align*}
		|K_\sigma(\xi)-K_\kappa(\xi)| & \leqslant \eps\int_{-1+\frac{|\xi|}{\gamma(b)}}^{-1+\frac{|\xi|}{\gamma(a)}} \frac{dy}{y+1} \\
		& \leqslant \eps \ln\left(\frac{\gamma(b)}{\gamma(a)}\right)\ .
	\end{align*}
	This concludes the first part of the proof. Now for such $\xi_0$ fixed lets take $\tilde{\delta}(\eps):=\frac{\delta(\eps)\gamma(a)}{\xi_0(\eps)}$ and $\kappa$ satisfying the hypothesis of the proposition. Then for any $|\xi|\leqslant \xi_0$ and $y\in[-1+\frac{|\xi|}{\gamma(b)},-1+\frac{|\xi|}{\gamma(a)}]$ we have
	\begin{align*}
		\left|\frac{y}{\sigma\circ\gamma^{-1}(\frac{|\xi|}{y+1})}- \frac{y}{\kappa\circ\gamma^{-1}(\frac{|\xi|}{y+1})}\right| &= |y|\frac{|\kappa\circ\gamma^{-1}(\frac{|\xi|}{y+1})-\sigma\circ\gamma^{-1}(\frac{|\xi|}{y+1})|}{\kappa\circ\gamma^{-1}(\frac{|\xi|}{y+1})\sigma\circ\gamma^{-1}(\frac{|\xi|}{y+1})} \\
		& \leqslant |y|\tfrac{\delta(\eps)\gamma(a)}{\xi_0} \\
		& \leqslant \left|-1+\tfrac{|\xi|}{\gamma(a)}\right|\,\tfrac{\delta(\eps)\gamma(a)}{\xi_0} \\
		& \leqslant \delta(\eps).
	\end{align*}
	Hence for $|\xi|\leqslant \xi_0$ (and the assumptions on $\kappa$ made in Proposition ~\ref{prop_continuity_kernel}) we obtain by uniform continuity of $\hat{h}$,
	\begin{equation*}
		|K_\sigma(\xi)-K_\kappa(\xi)|\leqslant \eps \ln\left(\frac{\gamma(b)}{\gamma(a)}\right).
	\end{equation*}
\end{proof}

The $|\hat{h}|^2$ uniform continuity hypothesis is neither absurd nor a constraint. Indeed if $h$ is in $L^1$ then by Riemann-Lebesgue lemma and property of the  Fourier transform we have that $\hat{h}\in C_0$ and so is $|\hat{h}|^2$. Having that leads to the uniform continuity of $|\hat{h}|^2$.

The following Lemma~\ref{lemma_up_bellow_control} results from the lower and upper bounds of the triangular inequality.

\begin{lemma}\label{lemma_up_bellow_control}
	Let $\sigma\in C^+([a,b])+1$ and $\kappa\in\mathcal{E}^+([a,b])+1$, then for any $\xi\in\R$,
	\begin{equation}
		|K_\kappa(\xi)|-|K_\sigma(\xi)-K_\kappa(\xi)| \leqslant |K_\sigma(\xi)| \leqslant |K_\sigma(\xi)-K_\kappa(\xi)| + |K_\kappa(\xi)|.
	\end{equation}
\end{lemma}

Putting together Lemma~\ref{lemma_up_bellow_control}, Proposition~\ref{prop_continuity_kernel}, Proposition~\ref{prop_lower_bound_kernel_basic} and Proposition~\ref{prop_control_norm_kernel_on_stair_function} yields Theorem ~\ref{theorem_norm_control_frequency_focus}.


